% META: {"id": "TSPE-CONTIN-EX-006", "chapitre": "TSPE-CONTINUITE", "type_objet": "exercice", "capacites": ["TSPE-CONTIN-C2"], "capacites_codes": ["C2"], "methodes": ["M2"], "parcours": 1, "competences": ["raisonner", "calculer"], "duree_min": 12, "mode_creation": "ex_nihilo", "sources_inspiration": [], "fichier_tex": "chapitres/TSPE-CONTINUITE/exercices/TSPE-CONTIN-EX-006.tex", "corrige_tex": "chapitres/TSPE-CONTINUITE/corriges/TSPE-CONTIN-CO-006.tex", "status": "approved"}
% BEGIN-VERIFY
% from sympy import symbols, sqrt, solve, Eq
% x = symbols('x')
% f = sqrt(x + 2)
% assert f.subs(x, 0) == sqrt(2)
% assert f.subs(x, 2) == 2
% sols = solve(Eq(x, f), x)
% assert sols == [2]
% END-VERIFY
\begin{exercice}{TSPE-CONTIN-EX-006}{1}{12}
Soit $(u_n)$ la suite définie par $u_0 = 0$ et $u_{n+1} = \sqrt{u_n + 2}$.
\begin{enumerate}
  \item Justifier que la fonction $f(x) = \sqrt{x+2}$ est continue sur $[0, 2]$ et que $f([0,2]) \subset [0,2]$.
  \item On admet que $(u_n)$ est croissante et majorée par $2$. Justifier qu'elle converge vers un réel $\ell \in [0, 2]$.
  \item Montrer que $\ell$ vérifie $\ell^2 - \ell - 2 = 0$, puis déterminer $\ell$.
\end{enumerate}
\end{exercice}
